\reslheading{经验、专长与技能}
  \begin{itemize}[leftmargin=*,label={}]
    \item \textbf{求真品诚，担当开发主力}
    {\small
      \begin{itemize}
        \item{基于LLVM的编译器开发，指令选择，DAG优化，别名分析，活跃变量分析等…}
        \item{为设计、验证提供可靠SDK环境；为算法提供高效驱动库；适量承担产品化任务。}
        \item{工程经验扎实，编译技术过硬：组织变革后，以一个月时间搭建tvm-llvm编译框架并跑通demo。}
      \end{itemize}
    }
    \item \textbf{响应积极，关注投入产出比}
    {\small
      \begin{itemize}
        \item{Pocket3产品手板阶段，SoC带宽争抢严峻，芯片对外DDR访存量遭质疑；硬件缺乏对应profiler，而仿真环境难以对齐产品板；利用编译插桩技术，直接在产品板上统计得到对外load/store指令量化数据，解决争议，反响极佳。}
        \item{芯片转阶段benchmark算子性能未达标；算法、设计已联合软硬件迭代无果；应用PGO、LTO技术，以较少投入大幅提升算子性能20\%，达标新一代芯片性能优化承诺，为项目阶段推进保驾护航。}
      \end{itemize}
    }
    \item \textbf{融会贯通，掌握芯片全景视角}
    {\small
      \begin{itemize}
        \item{负责自研芯片功耗拟合工作：产品化团队要求快速评估功耗；本资源组空有仿真profile数据，却未能聚合数据为简洁公式，仅将算子粗略分为计算密集型与访存密集型；提出使用Roofline Model拟合profile数据，以计算访存比为唯一输入，勾勒出自研处理器的性能功耗曲线。}
      \end{itemize}
    }
    \item \textbf{流利英语}
    {\small
      \begin{itemize}
        \item{曾负责对DJI瑞典哈苏团队的自研芯片培训。}
      \end{itemize}
    }
  \end{itemize} 